\originalchapter{磁场的来源}\label{ch:magnetostatics}
\originalsection{磁场与电流}
指南针转动说明空间存在能对磁针施加力矩的作用. 磁场$\bb$以特斯拉(T)为单位, 它对运动点电荷的力是$q\vct v\times\bb$. 磁场可以由电流、磁化材料及变化电场产生; 静止电荷本身在它的静止系中不产生磁场. 磁场线闭合或伸向无穷远, 没有像单个电荷那样已知的磁单极源. 这用$\oint\bb\cdot\dd\vct A=0$表示, 不意味着局部磁场为零.

稳恒电流元$I\dd\vct\ell$在观测点产生毕奥--萨伐尔场:
\[
\dd\bb=\frac{\mu_0}{4\pi}\frac{I\dd\vct\ell\times\uv R}{R^2},\qquad
\vct R=\rr-\rr'.
\]
源位置是$\rr'$, $\dd\vct\ell$沿电流方向. 该式属于稳恒电流的实验定律, 快速变化的电流不能简单把瞬时电流代入而忽略传播延迟. $\mu_0\simeq4\pi\times10^{-7}\,\mathrm{H/m}$, 此处写近似值而非现行SI精确常数.

\originalsection{直线、圆环与螺线管}
无限长直线沿$z$轴有电流$I$. 在轴距$s$处, 源点$z'$的距离平方为$s^2+z'^2$, 场方向是$\uv\phi$, 于是
\[
B_\phi=\frac{\mu_0I}{4\pi}\int_{-\infty}^{\infty}\frac{s\dd z'}{(s^2+z'^2)^{3/2}}
=\frac{\mu_0I}{2\pi s}.
\]
原函数为$z'/(s\sqrt{s^2+z'^2})$, 两端各给出$\pm1/s$. 电流向正$z$时, 从正$z$看场沿逆时针.

半径$R$圆环电流$I$, 轴上距中心$z$处横向分量对称相消. 电流元叉积的轴向因子为$R/\sqrt{R^2+z^2}$, 所以
\[
B_z=\frac{\mu_0I}{4\pi}\frac{R}{(R^2+z^2)^{3/2}}\int\dd\ell
=\frac{\mu_0IR^2}{2(R^2+z^2)^{3/2}}.
\]
中心场是$\mu_0I/(2R)$, 多匝同轴线圈在线圈厚度可忽略时乘匝数$N$.

\originalsection{安培定律与环路选择}
稳恒电流满足
\begin{equation}\label{eq:ampere}
\oint_C\bb\cdot\dd\vct\ell=\mu_0I_{\rm enc}.
\end{equation}
电流的正负依环路方向用右手法则定义. 它关联闭合线的环量与所围表面的穿越电流, 高斯定律则关联闭合面的通量与体内电荷. 两者的积分对象不可互换.

对长直线选同轴圆, 场沿切向且处处等大, 得$B2\pi s=\mu_0I$. 若电流填满半径$a$的均匀实心圆柱, 圆内电流为$I s^2/a^2$, 所以内场$B=\mu_0Is/(2\pi a^2)$、外场$B=\mu_0I/(2\pi s)$. 安培定律本身不保证对任意非对称电流只用一个环路就能求局部场.

理想无限长螺线管每单位长度有$n$匝. 内场轴向, 外场在无限长理想模型中为零. 穿过管内外的矩形环路长$\ell$, 有$B\ell=\mu_0n\ell I$, 故$B=\mu_0nI$. 有限管只在内部远离端部近似成立. 环形线圈有$N$匝, 在半径$s$的环形内部选圆, 得$B(s)=\mu_0NI/(2\pi s)$; 薄环近似才可视其为恒定.

无限薄电流面以面电流密度$K\uv z$沿$z$方向导电, $K$单位A/m. 跨面矩形环路给出两侧$B=\mu_0K/2$, 方向相反并由右手法则决定. 厚电流板可拆成许多这样的薄面, 在对称中心磁场相消; 不均匀$J(y)$必须对实际电流密度积分, 不能以总厚度比例代替.


安培圆不包围电流时环量为零, 也不说明局部场零. 例如在一根长直导线旁边取一个不围住导线的小矩形, 导线场在矩形每点仍非零, 只是各边环量相消. 对关于中心面偶对称的厚板电流$J_z(y)$, 没有额外背景场时可直接叠加薄面得到$B_x(y)=-\mu_0\int_0^yJ_z(u)\dd u$, 板内外都可用实际非零密度区积分; 第17章完整算非均匀例.

两同轴圆环中心关于原点对称, 同向电流产生同向中心场, 两个场的一阶轴向导数相消. Helmholtz线圈进一步把中心间距选为单环半径, 令二阶变化也相消, 在中心附近形成较均匀场. 相反电流则中心场抵消而梯度叠加, 适合展示非均匀场受力; “中心零场”并不表示周围无场.
